% !TEX program = pdflatex
% LLMWorld - extended paper (v3). pdfLaTeX-compatible, all-ASCII.
\documentclass[11pt]{article}

\usepackage[utf8]{inputenc}
\usepackage[T1]{fontenc}
\usepackage{lmodern}
\usepackage[a4paper,margin=2.4cm]{geometry}
\usepackage{amsmath}
\usepackage{amssymb}
\usepackage{array}
\usepackage{graphicx}
\usepackage{booktabs}
\usepackage{enumitem}
\usepackage[hidelinks]{hyperref}

\setlist{nosep,leftmargin=1.5em}

\title{\bfseries From Language to Load:\\ How Household LLM Agents Respond to Pricing, Norms and Disruptions\\ --- and Where the Response Stops Being Real}
\author{}
\date{September 2026}

\begin{document}
\maketitle

\begin{abstract}
\noindent
Residential electricity demand is increasingly shaped by heterogeneous household behaviour and by
distributed resources, yet the two dominant modelling traditions remain limited: top-down forecasting
treats households as an undifferentiated mass, and classical bottom-up models encode behaviour as fixed
states and transition probabilities that cannot process rich, contextual or textual signals such as
dynamic tariffs, social norms or news of a heatwave. Large language model (LLM) generative agents offer a
far more expressive behavioural engine, but existing LLM-based energy simulators are almost exclusively
single-household or small-scale, and none has been systematically validated as a population-level
policy-intervention testbed against the empirical demand-response literature. This paper presents
\textsc{LLMWorld}, an end-to-end multi-agent framework in which each household is an autonomous LLM agent
whose day is planned from a structured profile and then mapped to minute-resolution, appliance-level load;
crucially, every external input --- time-of-use or critical-peak prices, taxes, rebates, social norms, and
news events --- is delivered only in natural language, so that aggregate behaviour is emergent rather than
encoded. We run a large, controlled battery of experiments on three heterogeneous synthetic communities
under four validity rules derived from failure (same-era baselines, $n\ge15$, cross-world replication, and
content-contrast), and we deliberately subject every conclusion to an adversarial audit. We obtain seven
statistically significant cross-world event results, including a previously unreported stay-home family
--- public holidays, work-from-home days and transport strikes --- that each raise a household's daily
energy by roughly one fifth to two thirds, and we characterise the evening peak as cooking-dominated and
therefore largely resistant to generic signals. We further show that time-of-use and critical-peak prices,
which are null as plain text, become significantly peak-shifting once the trade-off is made concrete
through an injected per-appliance peak/off-peak cost table; decomposing this effect, the money figures
alone already produce a significant valley shift ($+11.6\%$), while an explicit rescheduling authorization
adds a further significant peak reduction ($-8.1\%$). Finally, we show that peak shaving obeys a mechanism
bounded by appliance substitutability: a device-targeted intervention shaves the total peak only for
non-substitutable devices (air-conditioning) and not for substitutable ones (induction cooking, which the
agent relocates to the oven and microwave). An adversarial self-audit withdraws one previously claimed
peak-shaving result that does not survive a shared-timeline replication, shows that all headline
$p$-values survive Benjamini--Hochberg correction, and classifies each finding as robust, prompt-sensitive
or withdrawn. We conclude by arguing that the calibrated capability boundary of LLM-driven energy
simulation --- not any single effect size --- is the contribution that makes it trustworthy as a
policy-experiment instrument.
\end{abstract}

\noindent\textbf{Keywords:} LLM agents; residential load simulation; demand response; natural-language
intervention; emergent behaviour; peak shaving; validity.

%======================================================================
\section{Introduction}

\subsection{Motivation}
Residential electricity consumption is becoming both more variable and more consequential. Rooftop
photovoltaics, batteries, electric vehicles and heat pumps are turning passive customers into active
prosumers, while demand-side flexibility is increasingly relied upon to absorb the variability of
renewable generation and to defer network reinforcement. Traditional load forecasting is ill-equipped for
this world. Top-down statistical models treat households as a single aggregate and cannot represent the
behavioural mechanisms through which prices, incentives or social signals actually change consumption.
Bottom-up models \cite{capasso1994,widen2010,richardson2010} reconstruct load from household composition
and activity patterns, but typically reduce behaviour to finite states or Markov chains with fixed
transition probabilities. Such representations struggle with exactly the signals that modern
demand-response programmes use: a dynamic tariff that must be understood and weighed, a social-comparison
message, or the news of an impending heatwave. Even agent-based hybrids \cite{xu2026} rely on rule-based
agents whose response repertoire is bounded by what the modeller encoded in advance.

Large language model (LLM) generative agents offer a fundamentally more expressive behavioural engine.
The ``generative agents'' architecture of Park et al.\ \cite{park2023} demonstrated that LLM agents
equipped with a memory stream, retrieval, reflection and planning can produce plausible, long-horizon
behaviour, and subsequent work has scaled such agents to large social simulations \cite{piao2025} and
reduced their computational cost through hierarchical prompting \cite{koley2025}. In the energy domain,
LLM agents have been used to infer single-household consumption from demographics \cite{chetty2025}, to
synthesise private household energy datasets without sensitive data \cite{almashor2024}, and to translate
non-standardised natural-language requests into structured parameters for a home energy management system
\cite{michelon2025}. What has \emph{not} been established is whether a population of heterogeneous LLM
households can be used as a \emph{policy experiment}: whether, when a tariff, a norm or a disruption is
injected purely as text, the population produces aggregate behaviour that is consistent with the
empirical demand-response literature --- and, equally important, where it is not.

This distinction matters for practice. A simulator that merely produces plausible-looking load curves is
a demonstration; a simulator that can tell a network operator, in advance, which interventions will and
will not move a population --- and by how much the answer can be trusted --- is an instrument. The aim of
this paper is to build the latter, and to be candid about its limits.

\subsection{Research questions}
The work is organised around four questions. \textbf{RQ1} asks how a large population of heterogeneous
household agents can be generated and mapped to realistic, diverse, minute-resolution load. In this paper
we exercise the population machinery at a deliberately modest scale and leave distributional
population-scale validation to future work; the contribution here is the pipeline rather than the scale.
\textbf{RQ2} asks how agents respond, at the appliance-decision level, when prices, taxes, subsidies,
norms and news events are injected as natural language --- that is, whether the behavioural transmission
mechanism that makes LLM agents attractive actually operates. \textbf{RQ3} asks whether the aggregate
responses align with the empirical benchmarks of Section~\ref{sec:benchmarks}, and whether phenomena
beyond individual behaviour --- in particular heterogeneity and dose-response --- emerge. \textbf{RQ4}
concerns scalability and practicality; it is explicitly out of scope as a research question, but its
implications surface in the discussion as a boundary condition.

\subsection{Contributions}
Our contributions are fivefold.
\begin{enumerate}
\item \textbf{Pipeline.} An end-to-end multi-agent pipeline that spans behaviour planning,
appliance-level decision-making and minute-resolution load generation, in which every external input is
delivered as natural language and no behavioural rule is hard-coded.
\item \textbf{Experimental programme.} A large, controlled battery of experiments across three
heterogeneous synthetic communities, obeying four validity rules derived from concrete failures, which
reports \textbf{seven statistically significant cross-world event results}, a validated price-response
capability, and a catalogue of well-characterised null results.
\item \textbf{A mechanism for peak shaving.} A device-targeted intervention shaves the total peak only
when it is anchored jointly to a \emph{device} and a \emph{window} \emph{and} the device is not
substitutable. This yields the organising observation that \emph{energy saving, peak shaving and load
shifting are three distinct behaviours} that one intervention may or may not produce.
\item \textbf{A calibrated capability boundary.} A precise account of what such a platform can and cannot
establish, including a \emph{withdrawn} result and the demonstration that all headline effects survive a
multiple-comparison correction.
\item \textbf{Methodological protocol.} Four validity rules --- same-era baselines, a minimum sample size,
cross-world replication, and content-contrast --- that we recommend for LLM-agent simulation studies
generally.
\end{enumerate}

%======================================================================
\section{Related work}

\subsection{From top-down forecasting to bottom-up behavioural models}
Bottom-up residential load models reconstruct demand from household and appliance inventories and from
occupant activity patterns. Capasso et al.\ \cite{capasso1994} combined socio-economic characteristics,
Monte-Carlo sampling of appliance ownership and probabilistic usage functions to reproduce
distribution-level demand without historical aggregate data. Wid\'en and W\"ackelg\r{a}rd \cite{widen2010}
derived activity sequences from time-use surveys via non-homogeneous Markov chains, and Richardson et al.\
\cite{richardson2010} jointly modelled occupancy and appliance usage to produce one-minute-resolution
profiles that matched both the distributions and the temporal structure of measured data. These models
remain the reference for synthetic load generation, and the present work reuses their core physical idea
--- behaviour mapped to appliance-level load at high resolution --- but replaces the Markov behavioural
engine with an LLM. The motivation is precisely the rigidity of the Markov representation: behaviour is
reduced to a small state space and fixed transition probabilities, and households are largely independent,
so social contagion, contextual reasoning and responses to textual or complex signals are out of reach.
Agent-based hybrids such as Xu et al.\ \cite{xu2026} add signal--demand feedback, but their agents remain
rule-based.

\subsection{LLM-based generative agents and reasoning}
The generative-agent architecture of Park et al.\ \cite{park2023} established a general template for
behaviour simulation. Scaling it raises engineering questions of cost and coherence: SALM
\cite{koley2025} reduced token consumption through hierarchical prompting and sustained thousands of
steps with an attention-based memory, while AgentSociety \cite{piao2025} scaled to more than ten thousand
urban agents and demonstrated system-level emergent patterns such as more equitable resource
distribution under broadcast policies. At the population level, Chen et al.\ \cite{chen2025} studied
consensus-seeking among LLM agents and showed that agent number, personality traits and network topology
govern convergence. The behavioural coherence of such agents ultimately rests on reasoning techniques:
chain-of-thought prompting \cite{wei2022} decomposes goals into ordered steps; ReAct \cite{yao2023a}
interleaves reasoning with environment actions; Tree-of-Thoughts \cite{yao2023b} searches over alternative
reasoning paths with lookahead and backtracking; and self-consistency \cite{wang2023} samples multiple
paths and majority-votes to reduce the variance injected by a single anomalous generation. These
techniques are the methodological substrate that lets an agent plan a day and then decide, segment by
segment, which appliances to run.

In the energy domain, Chetty et al.\ \cite{chetty2025} integrated chain-of-thought reasoning with a
reflective memory module to infer single-household consumption from demographic and occupational inputs,
reporting a mean absolute error below $10\%$ on the Pecan Street dataset; Almashor et al.\
\cite{almashor2024} used private LLMs to ``dream'' and execute daily behaviours and then extracted
appliance states by string matching, synthesising bimodal load curves without sensitive data; and
Michelon et al.\ \cite{michelon2025} combined ReAct with few-shot prompting to translate non-standardised
natural-language HEMS requests into structured parameters with high accuracy. These works establish
feasibility but remain single-household or small-scale, and none provides a population-level intervention
testbed with empirical alignment. That is the gap this paper addresses.

\subsection{Empirical benchmarks for alignment}
\label{sec:benchmarks}
The empirical literature provides quantitative targets against which a simulated population can be judged,
and we use them as alignment anchors throughout.

\noindent\emph{Pricing and demand response.} Faruqui and Sergici's survey of fifteen dynamic-pricing
experiments \cite{faruqui2010} reports that plain time-of-use pricing cuts peak demand by roughly 3--6\%,
that combining it with enabling technologies raises the cut to 10--30\%, and that critical-peak pricing
reaches 13--20\%. Demand-charge management prototypes report 10--20\% peak reduction \cite{escarrega2025}.
Against these, Wang et al.\ \cite{wang2021}, using provincial panel data and questionnaire responses
analysed with two-stage least squares, argue that habit rather than price is the dominant determinant of
household consumption --- an argument we return to when interpreting our own price nulls.

\noindent\emph{Non-price behavioural interventions.} Social-comparison feedback savings are of the order
of 1--3\%, with the strongest response among high-usage households \cite{allcott2011,ayres2013}. Loss
framing saves about five percentage points more than gain framing \cite{ghesla2019}. Monetary incentives
can crowd out intrinsic, social motives \cite{pellerano2017}; the effects of nudges and rebates are
strongly heterogeneous across households, favouring targeted delivery \cite{murakami2022}; and the same
intervention can work in opposite directions across ideological groups \cite{costa2010}.

\noindent\emph{Information.} Real-time in-home feedback can triple the price elasticity of demand
\cite{jessoe2012}, and personalised feedback outperforms aggregate billing \cite{monacchi2015}.

\noindent\emph{Structural factors.} Household-size economies imply that per-capita demand rises as
household size falls \cite{schroder2013}; active and passive peer effects shape the spatial diffusion of
residential solar panels \cite{barnes2022}; and comfort perception is the main obstacle to heating
sobriety in collective housing \cite{cabezas2025}.

\noindent\emph{Disruptions.} Empirically grounding LLM agents improves the realism of their behaviour
during disruptions \cite{xia2026}, and counterfactual simulation provides a control paradigm in which the
same population is observed with and without an event \cite{fidone2025}.

\noindent\emph{Platforms, tooling and methods.} CoRenew \cite{zhang2026} and Eco3S \cite{wei2026} are
LLM-agent and agent-based policy-simulation platforms; Gunkel et al.\ \cite{gunkel2023} analyse the
redistributional consequences of a uniform electricity tax across household categories; Zhou et al.\
\cite{zhou2016} show, using machine learning on observational smart-meter data, that users whose
consumption is more variable are more likely to reduce demand during demand-response events; residential
load-shape clustering underpins the segmentation of consumers \cite{jin2021}; and non-intrusive load
monitoring and multi-appliance forecasting provide the disaggregation and prediction baselines that a
synthetic population should be able to match.

\subsection{Positioning}
Our framework sits at the intersection of these strands. Like bottom-up models, it produces
minute-resolution, appliance-level load; like LLM generative agents, it plans behaviour in natural
language; and unlike prior LLM energy work, it is explicitly constructed and stress-tested as a
population-level \emph{policy-intervention instrument} whose outputs are compared, qualitatively, against
the empirical benchmarks above --- with an emphasis on honesty about where the comparison fails. To our
knowledge this is the first LLM-agent energy testbed that reports a withdrawn result and a
multiple-comparison correction as part of its primary contribution rather than as an afterthought.

%======================================================================
\section{The \textsc{LLMWorld} framework}

\subsection{Architecture and design principles}
A single entry point, \texttt{run.py}, exposes two modes. The \emph{world} mode synthesises a community in
four stages: household types, persona alignment, household construction and world assembly. The
\emph{simulate} mode generates behaviour, also in four stages: macro planning, inter-member coordination,
contextual enrichment and appliance decision-making. Three design principles run through the framework.
First, \textbf{all interventions are natural language}: a policy or an event is rendered into the prompt
as text, and the effect, if any, must emerge from the agent. Second, \textbf{every intermediate artefact
is persisted} as structured JSON, so that any experiment can be reconstructed and re-analysed offline
without re-querying the model; this is what makes the adversarial audit of Section~\ref{sec:audit}
possible. Third, \textbf{the physics is separate from the behaviour}: a dedicated load model converts
appliance decisions into energy, so that a change in behaviour is never silently confounded with a change
in how energy is computed.

\subsection{Domain model}
A household is modelled as a composite of \emph{members}, \emph{rooms} and \emph{appliances}. Appliances
are abstracted into four behavioural classes that determine the valid action space:
\emph{on-demand} (power is drawn only while the device is used, e.g.\ a lamp or an air-conditioner),
\emph{charging} (storage--discharge, e.g.\ a phone or an electric vehicle), \emph{cycle} (a fixed program
with a fixed energy per run, e.g.\ a washing machine, dishwasher or oven), and \emph{always-on}
(continuous baseload, e.g.\ a refrigerator). A shared catalog supplies, for each supported appliance
type, a rated power and plausible bounds, always-on daily energy, and experiment metadata --- crucially a
\texttt{flexible} flag, a duty cycle and a season tag --- as well as battery capacity and state of charge
for charging appliances, which are additionally bounded by a battery deficit so that a device cannot
absorb more than one full battery per day even if the agent requests it.

\subsection{The behaviour-to-load pipeline}
\noindent\textbf{Macro planning.} Each member receives a structured profile (age, occupation, habits,
personality) and the day's context (date, weather, active news) and produces a complete 00:00--24:00
timeline of activity segments. The output is constrained to real rooms, must cover the whole day, and
segments must be adjacent; a cross-day carry-over mechanism injects the previous day's end state so that
an agent does not ``teleport home'' at midnight --- a failure mode we observed and fixed early.

\noindent\textbf{Coordination.} Where a household has several members, their timelines are reconciled so
that, for example, two members do not occupy the kitchen simultaneously.

\noindent\textbf{Enrichment.} Each activity is enriched with the operational detail required to decide
appliance usage --- for instance, ``cook dinner'' is expanded into the sequence of appliance-level
operations it implies.

\noindent\textbf{Appliance decision.} This is the stage at which interventions bite. For every time
segment of the enriched timeline, the agent decides which appliances to operate, choosing among the
actions valid for each appliance class. The output is validated against the household registry: unknown
identifiers are repaired by appliance family, and operations on room appliances are dropped when the
member is out of the house.

\noindent\textbf{Energy calculation.} A separate physical model converts the appliance decisions into a
minute-resolution household load profile (1440 values) and a per-appliance energy breakdown. Cycle
appliances consume their fixed cycle energy, on-demand appliances their rated power weighted by duty
cycle, and idle appliances their standby draw. The peak-window energy used throughout the paper is the
integral of the profile over 16:00--21:00, and the valley over 22:00--07:00.

\subsection{Intervention surface}
All interventions are injected as natural language; no behavioural rule is hard-coded. \emph{Pricing}
policies render tariff text into the decision prompt and include time-of-use pricing (with a facts-only
variant), critical-peak pricing, critical-peak rebates, a uniform electricity tax, off-peak subsidies, a
demand charge and an EV-delay incentive. \emph{Social} interventions include a fixed social-norm
comparison, its soft variant, a loss-framed norm and a dynamic peer comparison built from the previous
day's community mean. \emph{Behaviour guidance} includes night setback and a real-time in-home display.
Policies can be combined (e.g.\ \texttt{tou,nudge}) and scheduled over a date window, so that
announcement, effectiveness and removal can be studied. \emph{News events} comprise thirteen preset
templates --- heatwave, cold-snap, storm, price-rise, supply-margin warning, air-conditioner peak tax,
rebate, rolling-blackout warning, solar subsidy, lockdown, public holiday, work-from-home day and
transport strike --- plus arbitrary custom events and community notices. Preset environment events
(heatwave, cold-snap) additionally adjust the structured weather fields so that the injected text and the
agent's environment agree.

\subsection{Concrete cost context}
Individually, none of the pricing policies above has any effect (Section~\ref{sec:price}); the agents
possess no economic model with which to weigh a tariff. We therefore introduce an optional
\emph{cost-context}: given the active tariff and the household's flexible appliances, the decision prompt
receives a concrete money table. The following is an actual injected example:
\begin{verbatim}
Flexible-appliance costs under today's tariff (peak 16:00-21:00 @0.90 AUD/kWh;
off-peak 22:00-07:00 @0.18 AUD/kWh):
- bedroom_1_airconditioner: 1.80 AUD now vs 0.36 AUD off-peak (save 1.44 per hour)
- kitchen_dishwasher:      0.99 AUD now vs 0.20 AUD off-peak (save 0.79 per cycle)
- bathroom_waterheater:    2.70 AUD now vs 0.54 AUD off-peak (save 2.16 per hour)
\end{verbatim}
Two variants are supported. In the \emph{strong} variant an explicit sentence authorizes the agent to
move a flexible appliance to an off-peak segment; in the \emph{soft} variant the same figures are given
with a neutral instruction, providing a prompt-bias control. A demand-charge variant instead renders each
appliance's kW times rate contribution to the daily maximum, and a multi-day variant injects yesterday's
peak-window cost as feedback.

\subsection{Analysis tooling}
A dedicated analysis package computes the quantities used throughout the paper: aggregate population
profiles, per-appliance usage, behaviour--load consistency checks, grouped responses (by energy awareness
or conscientiousness), event response, policy trade-offs, cross-world comparison, daily load-shape
clustering, hour-ahead forecastability with a random-forest model, household-size economies, and a peak
decomposition that isolates a single appliance family's contribution to the peak window. All outputs are
persisted as JSON or CSV. An offline suite of 298 unit tests, with the language model mocked, protects the
pure logic (schema parsing, policy rendering, tariff arithmetic, cleaning and repair) against regression.

%======================================================================
\section{Experimental design and validity controls}

\subsection{Worlds and households}
We use three synthetic communities with deliberately different structures, summarised in
Table~\ref{tab:worlds}. They differ in the number of households, in household size, and --- importantly
--- in where the evening peak originates, so that a mechanism found in one is not trivially a mechanism
found in all. Unless stated otherwise, an experiment targets one household on one simulated day and is
repeated many times.

\begin{table}[ht]
\centering\small
\begin{tabular}{lll}
\toprule
World & Composition & Evening-peak character \\
\midrule
W1 (\texttt{world\_838587}) & 2 households (1 and 4 members) & cooking-led; one nocturnal AC \\
W2 (\texttt{world\_172148}) & 3 households (5, 2 and 5 members) & cooking-led; living-room AC \\
W3 (\texttt{world\_143345}) & 3 households (6, 2 and 3 members, with EV) & mixed; AC-led household \\
\bottomrule
\end{tabular}
\caption{The three synthetic communities. Household sizes range from one to six members.}
\label{tab:worlds}
\end{table}

\subsection{Design and statistics}
Each comparison is a paired design in which the baseline and the treatment are interleaved so that they
share the same session-time context (``same-era''). Where the intervention acts at the
appliance-decision stage --- pricing and cost-context --- we run only that stage against a \emph{shared}
planning timeline, which removes planning-level noise and isolates the decision-level effect. Where the
intervention changes \emph{what the household does} --- structural events --- we run the full pipeline,
because the relevant behaviour is decided during planning; using a shared timeline for such events
produces a spurious null, a lesson we learned directly. Sample sizes are $n=15$ per arm unless noted;
effect sizes are reported as paired differences with paired $t$ tests, and $p$-values are two-sided. For
binary outcomes we use Fisher's exact test.

\subsection{Four validity rules}
A substantial part of the engineering effort went into discovering --- and then enforcing --- four rules
without which the results are not interpretable.
\begin{enumerate}
\item \textbf{Drift.} The continuous behaviour of an LLM agent drifts systematically across session time;
two independently sampled runs can differ by more than a treatment effect. We therefore always compare
against a same-era baseline, and, wherever possible, prefer binary or structural signals (appliance
on/off; out-of-home yes/no), which are immune to drift. The discovery of this artefact caused us to
withdraw two earlier peak-shaving results.
\item \textbf{Power.} Continuous magnitudes are unreliable below $n\approx15$: in a time-of-use
experiment, the estimated peak change moved from $-17.5\%$ at $n=3$ to $+13.1\%$ at $n=9$ to $-3.2\%$
(non-significant) at $n=15$. The run-to-run noise floor is a standard deviation of about $1.0$\,kWh
($\mathrm{CV}\approx12\%$), so literature-scale effects of 3--6\% lie below the resolution of the
instrument.
\item \textbf{Cross-world.} A mechanistic claim is accepted only if it replicates across worlds; several
initially promising effects were specific to a single household.
\item \textbf{Content, not channel.} An event delivered through the preset template and the same event
delivered as a custom string --- with byte-identical text --- differ on no metric, so the operative
variable is the \emph{content} of the intervention, not the injection path. This resolved a confound that
had confused two earlier experiments.
\end{enumerate}

%======================================================================
\section{Results}
\label{sec:results}

\subsection{Environmental and disruption events}
Table~\ref{tab:events} summarises the environmental and disruption events. A heatwave warning causes the
air-conditioner to be switched on where it was previously off: across member-days the air-conditioner was
active in $0/9$ baseline cases and $11/12$ heatwave cases (Fisher two-sided $p\approx3\times10^{-5}$),
consistent with the disruption-robustness argument of Xia et al.\ \cite{xia2026}. A cold-snap event
switches heating from off to on ($0/10$ vs $6/10$, $p\approx0.011$), with a weaker activation probability
than the heatwave --- cooling and heating events are not symmetric. A lockdown removes almost all
out-of-home activity ($13/13\rightarrow1/13$, $p\approx3\times10^{-6}$) and raises daytime (09:00--17:00)
energy by $220.7\%$ in a same-era design. We deliberately do not claim the latter magnitude: $220.7\%$ is
implausibly large relative to real lockdowns, and we read it as evidence that the agents over-comply with
an explicit stay-at-home directive. Finally, a severe-storm event is null, and a supply-margin warning
and a solar-subsidy announcement are null --- pure information without a concrete action does not move
behaviour, a theme that recurs below.

\begin{table}[ht]
\centering\small
\begin{tabular}{llll}
\toprule
Event & Type & Evidence & Significance \\
\midrule
Heatwave & environment & baseline AC $0/9$ vs $11/12$ & $p\approx3\times10^{-5}$ \\
Lockdown & structural & out-of-home $13/13\rightarrow1/13$; daytime $+220.7\%$ & $p\approx3\times10^{-6}$ \\
Cold-snap & environment & heating $0/10$ vs $6/10$ & $p\approx0.011$ \\
Storm & information & no change & n.s. \\
Supply-margin warning & information & no change & n.s. \\
Solar subsidy & information & no change & n.s. \\
\bottomrule
\end{tabular}
\caption{Environmental, structural and informational events. Only the first three are significant; the
pure-information events are null.}
\label{tab:events}
\end{table}

\subsection{The stay-home family}
A second, previously unreported group of results concerns \emph{structural stay-home} shocks. A public
holiday raises total daily energy in all seven households tested, by between $19\%$ and $65\%$
($p\le0.019$). A work-from-home day raises it by $30.6\%$, $28.7\%$ and $30.6\%$ in the three worlds
($p\le0.014$), which is the most consistent effect we observe. A transport strike raises it by $33.9\%$,
$22.2\%$ and $29.3\%$ ($p\le0.039$). Table~\ref{tab:stayhome} collects these.

The common mechanism is \emph{time at home}. The magnitude of the effect scales with how much the
household is normally away during the day: the holiday effect correlates negatively with the baseline
daytime load share ($r=-0.47$), and the weekend-versus-weekday effect correlates strongly with the
household's holiday effect ($r=0.82$), so the weekend is, in effect, a periodic holiday. Mechanistically,
the day's load is shifted earlier --- morning and midday equipment use rises while the late-afternoon load
falls --- which is exactly what happens when a commuter stays home. The evening-peak direction, by
contrast, is household-dependent ($-15\%$ to $+232\%$) and we do not assert it. A full-household
validation with all members simulated confirms the effect is not an artefact of simulating a single
member (holiday $+44.9\%$, work-from-home $+36.8\%$, transport strike $+50.3\%$).

\begin{table}[ht]
\centering\small
\begin{tabular}{lll}
\toprule
Event & Total-energy change & Significance \\
\midrule
Public holiday & $+19\%$ to $+65\%$ (7 households) & $p\le0.019$ (7/7) \\
Work-from-home & $+30.6\%$ / $+28.7\%$ / $+30.6\%$ & $p\le0.014$ \\
Transport strike & $+33.9\%$ / $+22.2\%$ / $+29.3\%$ & $p\le0.039$ \\
Weekend (reference) & $+3\%$ to $+45\%$ & descriptive \\
\bottomrule
\end{tabular}
\caption{The stay-home family: structural shocks that raise total energy by increasing time at home. The
weekend is the periodic version of the same mechanism.}
\label{tab:stayhome}
\end{table}

\subsection{Pricing as plain text is null}
As plain text, none of the pricing policies has a resolvable effect. Time-of-use pricing changes the
evening peak by $-3.2\%$ at $n=15$ (n.s.), and at smaller samples the estimate is unstable (Section~\ref{sec:design}).
A 7-fold change in the peak/valley rate gap has no effect. Critical-peak pricing does not shave the peak.
Adding a real-time in-home display to time-of-use pricing does not help. A demand charge is null. A
uniform electricity tax is null. Table~\ref{tab:pricenull} collects these text-only nulls. The pattern is
not ``prices do not matter'' --- it is that a tariff is, to these agents, a sentence rather than an
incentive: there is no cost to compute and no specific action to take.

\begin{table}[ht]
\centering\small
\begin{tabular}{lll}
\toprule
Intervention (text only) & Outcome & Note \\
\midrule
Time-of-use pricing & peak $-3.2\%$ (n.s.) & n=15 \\
TOU rate-gap sensitivity (7$\times$) & no effect & dose-response absent \\
Critical-peak pricing & peak null & \\
TOU + in-home display & null & no $\times$3 amplification \\
Demand charge & null & global objective \\
Uniform electricity tax & null & abstract \\
\bottomrule
\end{tabular}
\caption{Pricing interventions delivered as plain text are null.}
\label{tab:pricenull}
\end{table}

\subsection{Concrete cost information turns prices on}
\label{sec:price}
The picture changes decisively once the cost context is injected. Web at first tested a general
instruction to schedule flexible appliances at their cheapest feasible time and observed no effect; the
decisive ingredient is a concrete per-appliance money table (Section~3.5). Relative to a no-policy
baseline, on a shared timeline, time-of-use pricing with the cost context shifts the peak by $-13.2\%$,
$-5.0\%$ and $-3.5\%$ across the three worlds and raises the valley by $+25.6\%$, $+8.2\%$ and $+8.8\%$
($p\le0.035$); critical-peak pricing with the cost context gives a peak change of $-12.2\%$ and a valley
change of $+23.2\%$; and the critical-peak rebate gives $-13.7\%$ and $+27.9\%$. Total energy is
essentially unchanged in all cases: this is load shifting, not conservation. Table~\ref{tab:priceon}
summarises these.

Crucially, we decompose the effect. When the money table is delivered \emph{without} the authorizing
sentence (the soft variant), the valley still rises significantly by $+11.6\%$ ($p=0.019$) while the peak
falls only non-significantly by $-5.5\%$; adding the explicit authorization then produces a further
significant peak reduction of $-8.1\%$ ($p=0.031$) and a valley rise of $+12.5\%$ ($p=0.028$). The
interpretation is therefore two-part: \emph{concrete cost information is itself effective}, and an
actionable permission strengthens the peak effect --- but neither is a clean economic elasticity. The
effect is also dose-dependent: a larger peak/valley gap increases shifting (peak $-7.3\%$, valley
$+14.8\%$, $p\le0.014$), and households whose personas are more price-sensitive shift significantly more
(peak $-7.0\%$, valley $+14.6\%$, $p=0.014$). Reward and penalty framing are statistically
indistinguishable (critical-peak rebate $\approx$ critical-peak tax), and adding a social-norm or
loss-framed message on top of the concrete price changes the peak by only $+1.4\%$ and $+4.7\%$
respectively, both non-significant: once the window and the amount are concrete, framing stops mattering.

\begin{table}[ht]
\centering\small
\begin{tabular}{lll}
\toprule
Comparison & Peak 16--21h & Valley 22--7h \\
\midrule
TOU + cost vs baseline (W1/W2/W3) & $-13.2\%$ / $-5.0\%$ / $-3.5\%$ & $+25.6\%$ / $+8.2\%$ / $+8.8\%$ \\
CPP + cost vs baseline & $-12.2\%$ & $+23.2\%$ \\
CPR + cost vs baseline & $-13.7\%$ & $+27.9\%$ \\
Facts-only cost vs baseline & $-5.5\%$ (n.s.) & $+11.6\%$ ($p=0.019$) \\
Authorized vs facts-only & $-8.1\%$ ($p=0.031$) & $+12.5\%$ ($p=0.028$) \\
Larger price gap & $-7.3\%$ & $+14.8\%$ \\
High- vs low-sensitivity persona & $-7.0\%$ & $+14.6\%$ \\
\bottomrule
\end{tabular}
\caption{Concrete cost information turns null prices into significant shifting. The money figures alone
already shift the valley; the authorization adds peak reduction.}
\label{tab:priceon}
\end{table}

\subsection{Remaining price nulls}
Three pricing results remain null and are informative. A demand charge stays null even with concrete
figures, because it requires a \emph{global} ``lower today's maximum'' action that the agents do not
perform --- they can move a load out of a window, but they do not stagger appliances to compress the daily
maximum. A ``free electricity'' window during midday raises midday load nominally by $90\%$ but is not
statistically significant, and even the strongest possible wording does not produce a resolvable shift.
Multi-day bill feedback is a clean null: showing the household its previous day's bill does not change the
next day's behaviour, indicating an absence of cross-day memory or learning.

\subsection{The peak-shaving mechanism}
The evening peak is dominated by cooking in most households (the induction cooker accounts for
27--54\% of the peak-window energy), and this determines what can be shaved. The governing rule is that
targeting a device shaves the \emph{total} peak only if the device is not substitutable.
\emph{Targeting the air-conditioner}, which has no substitute, reduces air-conditioner energy by $58\%$
(device level). \emph{Targeting the induction cooker}, which does, reduces its own load by $99.7\%$ but
the agent simply relocates the activity to the oven and microwave, leaving the total peak unchanged
($-0.2\%$); a cooker peak tax is likewise null. \emph{Targeting the water heater} is ineffective in three
separate tests. Table~\ref{tab:shave} collects these. A device-targeted intervention therefore requires a
joint \emph{device} $\times$ \emph{window} anchor, and generic window-only signals do not shave.

This yields the organising observation that \emph{energy saving, peak shaving and load shifting are three
distinct behaviours}: a request can cut air-conditioner use while leaving the peak unmoved (saving
without shaving); a critical-peak price can shift load without saving (shifting without saving); and only
a device-and-window intervention on a non-substitutable device produces true peak shaving. The
device-targeted dose law --- shaving is proportional to the device's peak-window load times the achieved
behaviour-change rate --- holds as a trend across the three worlds, but it is bounded above by
substitutability.

\begin{table}[ht]
\centering\small
\begin{tabular}{llll}
\toprule
Target device & Substitutable? & Device effect & Total-peak effect \\
\midrule
Air-conditioner & no & $-58\%$ & shaved (device-level) \\
Induction cooker & yes & $-99.7\%$ & $-0.2\%$ (relocated to oven) \\
Water heater & (deferrable) & no response & null \\
\bottomrule
\end{tabular}
\caption{Peak shaving is bounded by substitutability: only non-substitutable devices yield total-peak
reduction.}
\label{tab:shave}
\end{table}

\subsection{Behavioural and social interventions}
Consistent with the pricing nulls, the non-price behavioural interventions are weak or unstable. Fixed
social comparison is directionally inconsistent --- it sometimes increases consumption --- and a
loss-framed norm, which the literature would predict to be stronger, produces a counterexample (a
$+6\%$ change in the wrong direction). In-home display and night setback are null. Two combinations are
informative. Time-of-use pricing combined with a social norm is non-additive: the measured effect of
$-13.9\%$ is far from the additive prediction of $-2.3\%$. And the air-conditioner peak tax combined with
a generic rebate \emph{reverses} the tax's effect: the tax alone reduces air-conditioner energy by $58\%$,
but adding a rebate brings it back up by $80\%$, consistent with the crowding-out documented by Pellerano
et al.\ \cite{pellerano2017}. Table~\ref{tab:social} collects the behavioural results.

\begin{table}[ht]
\centering\small
\begin{tabular}{lll}
\toprule
Intervention & Outcome & Reference \\
\midrule
Social norm (nudge) & unstable direction & Allcott 2011; Ayres 2013 \\
Loss-framed norm & counterexample ($+6\%$) & Ghesla 2019 \\
Peer comparison (dynamic) & inconsistent & Ayres 2013 \\
In-home display & null & Jessoe \& Rapson 2012 \\
Night setback & null & Cabezas-Riviere 2025 \\
TOU + nudge & non-additive & Pellerano 2017 \\
AC tax + rebate & crowding-out & Pellerano 2017 \\
Framing on top of concrete price & no effect & --- \\
\bottomrule
\end{tabular}
\caption{Behavioural and social interventions.}
\label{tab:social}
\end{table}

\subsection{Descriptive analyses of the synthetic population}
Several analyses characterise the generated population independently of any intervention. Households split
into recognisable load-shape archetypes --- evening-peak, morning-peak and late-night types --- mirroring
the segmentation reported for real smart-meter data \cite{jin2021}. Per-capita daily energy falls as
household size rises, qualitatively consistent with the household-size economies of Schr\"oder et al.\
\cite{schroder2013}. Weekend energy exceeds weekday energy by $3\%$ to $45\%$ depending on the household,
the periodic version of the stay-home effect. Season-correct behaviour emerges once the season is derived
from the date: winter households use heating and no cooling, summer households the reverse. A
random-forest one-hour-ahead forecaster improves on a persistence baseline by $31.5\%$ to $57.1\%$,
indicating that the generated load is structurally regular rather than noise. Finally, households whose
baseline consumption is more variable tend to respond more strongly to pricing; the correlation is in the
direction reported by Zhou et al.\ \cite{zhou2016} but is not statistically significant at the five
households available ($r=0.51$).

\subsection{Summary of findings}
Table~\ref{tab:summary} summarises the principal findings and their status after the audit of
Section~\ref{sec:audit}.

\begin{table}[ht]
\centering\small
\begin{tabular}{lll}
\toprule
Finding & Result & Status \\
\midrule
Heatwave, cold-snap, lockdown & significant, cross-world & robust \\
Stay-home family (holiday/wfh/strike) & $+19\%$ to $+65\%$, cross-world & robust \\
Concrete-priced TOU/CPP/CPR & peak $-3.5\%$ to $-13.2\%$, valley $+8\%$ to $+28\%$ & robust \\
Facts-only vs authorized cost & valley $+11.6\%$; peak $-8.1\%$ & robust (decomposed) \\
Device-and-window shaving (AC) & device $-58\%$ & robust (device level) \\
Substitution (cooker) & no total-peak shaving & robust \\
Text-only prices, demand charge & null & robust \\
Multi-day feedback, free window & null & robust \\
AC peak-tax total-peak $-22.5\%$ & withdrawn (shared timeline $+6.1\%$) & withdrawn \\
AC-tax $\times$ rebate (n=9) & non-additive & weak \\
Variability-to-response (n=5) & direction only & weak \\
\bottomrule
\end{tabular}
\caption{Summary of principal findings and their audited status.}
\label{tab:summary}
\end{table}

%======================================================================
\section{Discussion}
\label{sec:discussion}

\subsection{Two classes of behaviour}
The clearest reading of our results is that the platform contains two qualitatively different kinds of
behaviour. Structural, contextual shocks --- weather and stay-at-home events --- change the situation the
agent is reasoning about, and the agent responds through ordinary commonsense reasoning; these effects
are large, robust and cross-world. Generic prices, by contrast, demand an economic trade-off that the
agents do not perform: a tariff is, to them, a sentence, not an incentive. They respond only when the
trade-off is made \emph{concrete} (a per-appliance money figure) and, for peak shaving, \emph{actionable}
(a device and a window). The decomposition in Section~\ref{sec:price} sharpens this: the money figures
alone are already sufficient for a significant valley shift, but not for a significant peak reduction, so
the peak effect specifically requires an actionable permission. It is tempting to summarise this as ``the
agents understand prices''; the more accurate statement is that they act on concrete, actionable
information, and that supplying the money figures is what makes the instruction concrete.

\subsection{Alignment with the empirical literature}
The \emph{directions} of our results align with the empirical record (Table~\ref{tab:align}). Peak
shifting under concrete pricing is consistent with Faruqui and Sergici \cite{faruqui2010}; the disruption
results with Xia et al.\ \cite{xia2026}; the non-additivity of combinations and the crowding-out of a
device tax by a rebate with Pellerano et al.\ \cite{pellerano2017}; the fragility of framing effects
without concrete anchors with the habit-dominance argument of Wang et al.\ \cite{wang2021}; the
heterogeneous price response with Costa and Kahn \cite{costa2010} and Murakami et al.\ \cite{murakami2022};
and the variability-to-response direction with Zhou et al.\ \cite{zhou2016}. The \emph{magnitudes},
however, systematically exceed the empirical values, which we attribute to LLM over-compliance with
explicit natural-language directives: an instruction-following agent overshoots relative to a habit-bound
human household. For this reason we treat alignment as qualitative and decline to claim absolute
magnitudes.

\begin{table}[ht]
\centering\small
\begin{tabular}{llll}
\toprule
Phenomenon & Empirical benchmark & Our result & Alignment \\
\midrule
Plain TOU & $-3$ to $-6\%$ & text null; cost-context $-3.5$ to $-13.2\%$ & qualitative \\
Critical-peak pricing & $-13$ to $-20\%$ & $-12.2\%$ with cost context; null as text & partial \\
Social norm & $1$ to $3\%$ & unstable & not aligned \\
Loss framing & $+5$ pp vs gain & counterexample & not aligned \\
Monetary vs norm & crowding-out & AC tax reversed by rebate & aligned \\
Disruption realism & behaviour shifts & significant, cross-world & aligned \\
Variability to DR & higher variability responds more & direction ($r=0.51$, n=5) & weak \\
\bottomrule
\end{tabular}
\caption{Alignment of our findings with the empirical benchmarks of Section~\ref{sec:benchmarks}.}
\label{tab:align}
\end{table}

\subsection{A calibrated capability boundary}
We therefore state the boundary explicitly. The platform \emph{can} reliably give: (i) the direction of
structural weather and stay-home events; (ii) peak shifting and valley filling under concrete pricing;
and (iii) the device-level response of non-substitutable appliances. It \emph{cannot} give
literature-scale (3--6\%) text-price elasticity, global demand-charge shaving, cross-day learning or
habit, and it cannot represent photovoltaic-adoption or electric-vehicle valley phenomena because those
mechanisms are absent. Naming this boundary is not a concession; it is what turns an impressive demo into
a usable instrument, because it tells a policy analyst in advance which questions the tool can answer and
where a result would be an artefact.

\subsection{Methodological contributions}
Three methodological findings may outlast the platform itself. First, continuous LLM-agent behaviour
drifts across session time, so same-era baselines or binary signals are mandatory --- a result that
invalidated two of our own earlier conclusions. Second, magnitudes below $n\approx15$ are unreliable and
should not be reported. Third, the effect of an intervention depends on its \emph{content}, not on the
channel through which it is delivered. Together with the four rules of Section~\ref{sec:design}, these
findings form a validity protocol that we recommend for LLM-agent simulation studies generally, and which
we would have liked to have had at the start of this project.

\subsection{Implications}
For practice, the most useful result may be the decomposition of price response into
information and authorization: it suggests that a tariff communicates effectively only when rendered as a
concrete, per-appliance trade-off, and that pure framing devices add little once that is done. For the
research community, the most useful result may be the boundary: an LLM-agent population is a powerful
generator of \emph{qualitative, directional} hypotheses about demand-side behaviour, but it is not yet a
quantitative forecaster, and treating it as one would be a mistake.

%======================================================================
\section{Adversarial validity audit}
\label{sec:audit}
Because the experimental programme is large, we subjected every conclusion to a deliberate attempt at
disconfirmation, and we report the outcome even where it is unflattering.

\noindent\textbf{Robust.} The seven event results and the null set are robust in direction. To guard
against multiple-comparison artefacts, we applied the Benjamini--Hochberg procedure across the eight
primary tests; all eight survive (the largest adjusted $p$-value is $0.039$). The device-level
air-conditioner reduction and the substitution result are robust because they are device-level effects
that do not depend on the evening-peak metric.

\noindent\textbf{Prompt-sensitive (preliminary).} The price-shaving result is partly instruction-driven
and should not be reported as a pure elasticity; the accurate claim is that concrete cost information is
effective and an actionable permission strengthens peak shaving. The heatwave result is not purely
language-only, because the preset event also adjusts a structured weather field, so it demonstrates
cross-modal consistency rather than language-only emergence. The season result largely follows the
prompt's own season rule and is therefore weak evidence of emergent seasonal behaviour.

\noindent\textbf{Withdrawn or demoted.} One previously claimed result --- that an air-conditioner peak tax
reduces the total evening peak by $22.5\%$ --- is \emph{withdrawn}: on a shared timeline the corresponding
comparison gives $+6.1\%$ (n.s.), and only the device-level air-conditioner reduction of $58\%$ survives.
Two further results are demoted to weak evidence: the non-additivity of the air-conditioner tax and rebate
is based on $n=9$, and the substitution-based explanation of the cooking peak did not replicate on a
full five-member household.

\noindent\textbf{Cross-cutting risks.} We also record risks that we could only partially mitigate:
over-compliance inflating magnitudes; interventions that touch structured fields rather than language
alone; a household-specific ``generic-event reactivity'' that inflates load on one community and confounds
several comparisons there; the post-hoc tuning of the authorization wording in the cost context; the
single-member scope of most experiments; and platform-version drift during the study, as analysis and
weather modules were extended mid-programme.

%======================================================================
\section{Limitations}
The weather model is a stub beyond the date-derived season and the event overrides, and the seasonal
temperatures are nominal. We did not align against public occupant-driven stochastic load data, so the
alignment with the empirical literature is qualitative rather than distributional. The population is
exercised at no more than five households per world, and the two research questions that concern scale and
cost are out of scope. Most experiments simulate a single household member, and multi-member validation
was performed only for a subset of effects. Absolute magnitudes are inflated by over-compliance and are
not claimed. Several phenomena of practical interest --- photovoltaic adoption, electric-vehicle valley
filling, and any form of cross-day learning --- cannot currently be studied on this platform. Finally, the
agents do not carry a memory across days, which limits the study of habit formation, a mechanism that the
empirical literature identifies as central.

%======================================================================
\section{Conclusion}
This paper set out to ask whether a population of heterogeneous LLM household agents can be used as an
instrument for natural-language energy-policy experiments, and to answer that question honestly. The
instrument we built, \textsc{LLMWorld}, produces emergent, minute-resolution behaviour across an
end-to-end behaviour-to-load pipeline. It yields seven statistically significant, cross-world event
results --- including the stay-home family --- a validated price-response capability whose effect
decomposes into concrete information and actionable authorization, and a mechanistic account of peak
shaving governed by a dose, substitutability and device/window rule. Equally important, it yields a
precise negative space: text-only prices, global demand charges, cross-day learning, and photovoltaic and
electric vehicle effects are not attainable, and one previously claimed result is withdrawn. We argue
that this calibrated boundary --- the explicit statement of what an LLM-agent energy simulator can and
cannot establish --- is the contribution that ultimately makes such platforms credible, and we release
the pipeline, the intervention surface and the full experimental record to support replication.

%======================================================================
\begin{thebibliography}{99}
\small
\bibitem{allcott2011} Allcott, H. (2011). Social norms and energy conservation. Journal of Public Economics, 95(9--10), 1082--1095.
\bibitem{almashor2024} Almashor, M., et al. (2024). Can private LLM agents synthesize household energy consumption data? Energy and AI, 16, 100360.
\bibitem{ayres2013} Ayres, I., Raseman, S., and Shih, A. (2013). Evidence from two large field experiments that peer comparison feedback can reduce residential energy usage. JLEO, 29(5), 992--1022.
\bibitem{barnes2022} Barnes, J., Islam, M. R., and Morshed, S. (2022). Passive and active peer effects in the spatial diffusion of residential solar panels. Energy Research and Social Science, 86, 102458.
\bibitem{cabezas2025} Cabezas-Riviere, N., et al. (2025). Identifying and understanding obstacles to heating sobriety and thermal comfort in collective housing. arXiv:2512.16949.
\bibitem{capasso1994} Capasso, A., Grattieri, W., Lamedica, R., and Prudenzi, A. (1994). A bottom-up approach to residential load modeling. IEEE Transactions on Power Systems, 9(2), 957--964.
\bibitem{chen2025} Chen, H., et al. (2025). Multi-agent consensus seeking via large language models. arXiv:2310.20151.
\bibitem{chetty2025} Chetty, N., et al. (2025). An LLM framework for inferring household energy consumption through behaviour simulation. Applied Energy, 355, 122335.
\bibitem{costa2010} Costa, D. L., and Kahn, M. E. (2010). Energy conservation nudges and environmentalist ideology. JEEA, 11(3), 680--702.
\bibitem{escarrega2025} Escarrega, C., et al. (2025). Demand charge management: Prototype design and testing. arXiv:2509.10713.
\bibitem{faruqui2010} Faruqui, A., and Sergici, S. (2010). Household response to dynamic pricing of electricity: A survey of 15 experiments. Journal of Regulatory Economics, 38(2), 193--225.
\bibitem{fidone2025} Fidone, G., et al. (2025). Evaluating online moderation via LLM-powered counterfactual simulations. arXiv:2511.07204.
\bibitem{ghesla2019} Ghesla, C., Grieder, M., and Schmitz, J. (2019). Pro-environmental incentives and loss aversion. Energy Policy, 128, 102--112.
\bibitem{gunkel2023} Gunkel, P. A., et al. (2023). Uniform taxation of electricity. arXiv:2306.11566.
\bibitem{jessoe2012} Jessoe, K., and Rapson, D. (2012). Knowledge is (less) power. American Economic Review, 104(4), 1417--1438.
\bibitem{jin2021} Jin, L., et al. (2021). Investigating underlying drivers of variability in residential energy usage patterns. arXiv:2102.11027.
\bibitem{koley2025} Koley, S. (2025). SALM: A multi-agent framework for language model-driven social network simulation. arXiv:2505.09081.
\bibitem{michelon2025} Michelon, F., Zhou, Y., and Morstyn, T. (2025). Large language model interface for home energy management systems. Applied Energy, 358, 122590.
\bibitem{monacchi2015} Monacchi, A., et al. (2015). An open solution to provide personalized feedback for building energy management. arXiv:1505.01311.
\bibitem{murakami2022} Murakami, K., Shimada, K., and Tanaka, M. (2022). Heterogeneous treatment effects of nudge and rebate. Energy Economics, 105, 105737.
\bibitem{park2023} Park, J. S., et al. (2023). Generative agents: Interactive simulacra of human behavior. Proc. UIST 2023.
\bibitem{pellerano2017} Pellerano, J. A., et al. (2017). Do extrinsic incentives undermine social norms? Environmental and Resource Economics, 67(3), 413--431.
\bibitem{piao2025} Piao, J., et al. (2025). AgentSociety: Large-scale LLM-driven generative agents for social simulation. arXiv.
\bibitem{richardson2010} Richardson, I., Thomson, M., Infield, D., and Clifford, C. (2010). Domestic electricity use: A high-resolution energy demand model. Energy and Buildings, 42(10), 1878--1887.
\bibitem{schroder2013} Schroeder, C., et al. (2013). Household formation and residential energy demand: Evidence from Japan. Energy Economics.
\bibitem{thorvaldsen2021} Thorvaldsen, K., et al. (2021). Long-term value of flexibility from flexible assets in building operation. arXiv:2105.11952.
\bibitem{wang2021} Wang, Y., et al. (2021). Electricity price and habits. Energy Policy.
\bibitem{wang2023} Wang, X., et al. (2023). Self-consistency improves chain of thought reasoning in language models. Proc. ICLR.
\bibitem{wei2022} Wei, J., et al. (2022). Chain-of-thought prompting elicits reasoning in large language models. Proc. NeurIPS.
\bibitem{wei2026} Wei, Y., et al. (2026). Eco3S: Complex socio-economic system simulation via agent-based models. arXiv:2607.26588.
\bibitem{widen2010} Widen, J., and Wackelgard, E. (2010). A high-resolution stochastic model of domestic activity patterns and electricity demand. Applied Energy, 87(6), 1880--1892.
\bibitem{xia2026} Xia, H., et al. (2026). Empirical grounding improves the realism of LLM agents simulating human behavior during disruptions. arXiv:2607.17437.
\bibitem{xu2026} Xu, et al. (2026). Agent-based modeling and neural network for residential customer demand response.
\bibitem{yao2023a} Yao, S., et al. (2023a). ReAct: Synergizing reasoning and acting in language models. Proc. ICLR.
\bibitem{yao2023b} Yao, S., et al. (2023b). Tree of thoughts: Deliberate problem solving with large language models. Proc. NeurIPS.
\bibitem{zhang2026} Zhang, et al. (2026). CoRenew: A large language model agent-based policy simulation platform for multifamily residential redevelopment. arXiv:2607.25447.
\bibitem{zhou2016} Zhou, D., Balandat, M., and Tomlin, C. (2016). Residential demand response targeting using machine learning with observational data. arXiv:1607.00595.
\end{thebibliography}

\appendix
\section{Reproducibility}
The repository contains the simulation entry point, the engine modules for pricing, cost-context,
news/events and weather, the world- and simulate-stage steps, the analysis package, 298 offline unit tests
with the language model mocked, and the complete experimental record of 222 rounds. Every intervention is
delivered in natural language; no behavioural rule is hard-coded.

\end{document}
